\documentclass[12pt,a4paper]{article}

\usepackage[utf8]{inputenc}
\usepackage[T1]{fontenc}
\usepackage[spanish]{babel}
\usepackage{geometry}
\usepackage{graphicx}
\usepackage{xcolor}
\usepackage{titlesec}
\usepackage{fancyhdr}
\usepackage{enumitem}
\usepackage{booktabs}
\usepackage{array}
\usepackage{tabularx}
\usepackage{listings}
\usepackage{hyperref}
\usepackage{mdframed}
\usepackage{tcolorbox}
\usepackage{amsmath}
\usepackage{amssymb}
\usepackage{parskip}
\usepackage{setspace}

\geometry{
  top=2.5cm, bottom=2.5cm,
  left=3cm, right=3cm,
  headheight=15pt
}

% Colors
\definecolor{accent}{RGB}{255, 92, 0}
\definecolor{navy}{RGB}{0, 27, 61}
\definecolor{surface}{RGB}{16, 20, 21}
\definecolor{codeblue}{RGB}{100, 160, 242}
\definecolor{codegray}{RGB}{140, 144, 153}
\definecolor{codesurface}{RGB}{29, 32, 34}
\definecolor{lightgray}{RGB}{245, 245, 245}

% Hyperref setup
\hypersetup{
  colorlinks=true,
  linkcolor=accent,
  urlcolor=accent,
  citecolor=navy,
  pdfauthor={Pick \& Serve Team},
  pdftitle={Pick \& Serve – White Paper Técnico},
}

% Section formatting
\titleformat{\section}
  {\large\bfseries\color{navy}}
  {\thesection.}{0.8em}{}[{\color{accent}\titlerule[1.5pt]}]

\titleformat{\subsection}
  {\normalsize\bfseries\color{navy}}
  {\thesubsection.}{0.6em}{}

\titleformat{\subsubsection}
  {\normalsize\bfseries\color{accent}}
  {\thesubsubsection.}{0.6em}{}

\titlespacing{\section}{0pt}{18pt}{8pt}
\titlespacing{\subsection}{0pt}{12pt}{4pt}

% Header / Footer
\pagestyle{fancy}
\fancyhf{}
\fancyhead[L]{\small\textcolor{codegray}{Pick \& Serve — White Paper Técnico}}
\fancyhead[R]{\small\textcolor{codegray}{Ingeniería de Software II}}
\fancyfoot[C]{\small\textcolor{codegray}{\thepage}}
\renewcommand{\headrulewidth}{0.4pt}
\renewcommand{\headrule}{\hbox to\headwidth{\color{accent}\leaders\hrule height \headrulewidth\hfill}}

% Code listings
\lstset{
  backgroundcolor=\color{lightgray},
  basicstyle=\ttfamily\small,
  keywordstyle=\color{navy}\bfseries,
  commentstyle=\color{codegray}\itshape,
  stringstyle=\color{accent},
  breaklines=true,
  frame=single,
  framerule=0.4pt,
  rulecolor=\color{codegray},
  numbers=left,
  numberstyle=\tiny\color{codegray},
  numbersep=6pt,
  captionpos=b,
  showspaces=false,
  showtabs=false,
  tabsize=2,
}

% Tcolorbox styles
\tcbuselibrary{skins,breakable}

\newtcolorbox{infobox}[1][]{
  colback=navy!5!white,
  colframe=navy,
  fonttitle=\bfseries,
  title=#1,
  breakable,
}

\newtcolorbox{accentbox}[1][]{
  colback=accent!8!white,
  colframe=accent,
  fonttitle=\bfseries\color{white},
  coltitle=white,
  title=#1,
  breakable,
}

\newcommand{\code}[1]{\texttt{\small\color{navy}#1}}
\newcommand{\event}[1]{\texttt{\small\color{accent}#1}}

\begin{document}

% ──────────────────────────────────────────────────────────────────────────────
% PORTADA
% ──────────────────────────────────────────────────────────────────────────────
\begin{titlepage}
  \thispagestyle{empty}

  \vspace*{2cm}
  \begin{center}
    {\color{accent}\rule{\linewidth}{3pt}}
    \vspace{0.5cm}

    {\huge\bfseries\color{navy} Pick \& Serve}\\[0.4cm]
    {\Large\color{accent} Predictor de Partidos ATP}\\[0.3cm]
    {\large White Paper Técnico}

    \vspace{0.5cm}
    {\color{accent}\rule{\linewidth}{1pt}}

    \vspace{2cm}

    \begin{tcolorbox}[colback=navy!5, colframe=navy!30, width=0.85\linewidth]
      \centering
      \textbf{Resumen Ejecutivo}\\[0.5em]
      \small
      Pick \& Serve es una plataforma web de predicción de partidos de tenis ATP,
      construida sobre una arquitectura orientada a eventos con RabbitMQ como broker de mensajería.
      El sistema permite a múltiples usuarios competir en un ranking en tiempo real, realizando
      pronósticos sobre partidos de torneos ATP. El presente documento describe la arquitectura
      técnica, decisiones de diseño, stack tecnológico y flujos de datos del sistema.
    \end{tcolorbox}

    \vspace{2.5cm}

    \begin{tabular}{ll}
      \textbf{Materia:}     & Ingeniería de Software II \\[4pt]
      \textbf{Tipo:}        & Trabajo Práctico — Prueba de Concepto \\[4pt]
      \textbf{Tecnologías:} & FastAPI · React · RabbitMQ · PostgreSQL \\[4pt]
      \textbf{Versión:}     & 1.0.0 \\[4pt]
      \textbf{Fecha:}       & Junio 2026 \\
    \end{tabular}

    \vfill
    {\color{accent}\rule{\linewidth}{2pt}}
  \end{center}
\end{titlepage}

% ──────────────────────────────────────────────────────────────────────────────
% TABLA DE CONTENIDOS
% ──────────────────────────────────────────────────────────────────────────────
\tableofcontents
\newpage

% ──────────────────────────────────────────────────────────────────────────────
\section{Introducción y Contexto}
% ──────────────────────────────────────────────────────────────────────────────

\subsection{Descripción del Problema}

Los sistemas de predicción deportiva presentan desafíos técnicos particulares que los hacen
idóneos para estudiar arquitecturas de software modernas:

\begin{itemize}
  \item \textbf{Alta concurrencia}: múltiples usuarios enviando predicciones simultáneamente.
  \item \textbf{Consistencia de datos}: los puntajes deben calcularse correctamente aunque
        varios eventos ocurran en paralelo.
  \item \textbf{Actualización en tiempo cuasi-real}: el ranking debe reflejar los resultados
        apenas se cargan los datos, sin requerir recarga manual.
  \item \textbf{Desacoplamiento}: las acciones administrativas (cargar resultados, cerrar
        jornadas) no deben bloquear al usuario final.
\end{itemize}

\subsection{Propuesta de Valor}

Pick \& Serve aborda estos desafíos implementando una \textbf{arquitectura orientada a eventos
(Event-Driven Architecture, EDA)} donde cada acción relevante del sistema se traduce en un
evento publicado en un broker de mensajería, permitiendo que múltiples consumidores reaccionen
de forma independiente y asíncrona.

\subsection{Alcance del Sistema}

El sistema cubre el flujo completo:
\begin{enumerate}
  \item Autenticación de jugadores y administradores.
  \item Carga de pronósticos por parte de los jugadores sobre partidos abiertos.
  \item Carga de resultados por parte del administrador.
  \item Puntuación automática via workers asíncronos.
  \item Notificaciones a jugadores sobre resultados y cierres de jornada.
  \item Ranking global actualizado automáticamente.
\end{enumerate}

% ──────────────────────────────────────────────────────────────────────────────
\section{Arquitectura General del Sistema}
% ──────────────────────────────────────────────────────────────────────────────

\subsection{Visión de Alto Nivel}

El sistema se organiza en cinco capas principales que se comunican de forma desacoplada:

\begin{accentbox}[Capas del Sistema]
\begin{enumerate}
  \item \textbf{Frontend (React + TypeScript)} — Interfaz web responsive (SPA).
  \item \textbf{Backend API (FastAPI)} — Endpoints REST, lógica de negocio, publicación de eventos.
  \item \textbf{Broker de Mensajería (RabbitMQ)} — Exchange topic para routing de eventos.
  \item \textbf{Workers Asíncronos (Python)} — Consumidores especializados: scoring, notificaciones, ranking.
  \item \textbf{Persistencia (PostgreSQL)} — Base de datos relacional con SQLAlchemy ORM.
\end{enumerate}
\end{accentbox}

\subsection{Diagrama de Componentes}

\begin{center}
\begin{tabular}{|c|}
  \hline
  \textbf{Cliente Browser} \\
  React SPA (Vite + TypeScript) \\
  \hline
\end{tabular}
$\downarrow$ HTTP/REST
\begin{tabular}{|c|}
  \hline
  \textbf{API Server} \\
  FastAPI + Uvicorn \\
  \hline
\end{tabular}
$\rightarrow$ AMQP $\rightarrow$
\begin{tabular}{|c|}
  \hline
  \textbf{RabbitMQ} \\
  Exchange: topic \\
  \hline
\end{tabular}

\vspace{0.5cm}
\small
RabbitMQ distribuye eventos a tres colas paralelas:
\event{scoring.queue} · \event{notifications.queue} · \event{ranking.queue}
\end{center}

\subsection{Principios de Diseño Aplicados}

\begin{description}
  \item[Separación de responsabilidades (SoC)] Cada worker tiene una única responsabilidad bien definida.
  \item[Alta cohesión / Bajo acoplamiento] Los módulos interactúan únicamente via contratos de eventos (mensajes JSON).
  \item[Idempotencia] Los workers pueden procesar el mismo mensaje múltiples veces sin generar efectos adversos.
  \item[Resiliencia] El broker persiste mensajes; si un worker cae, los mensajes se encolan y procesan al recuperarse.
\end{description}

% ──────────────────────────────────────────────────────────────────────────────
\section{Arquitectura Orientada a Eventos}
% ──────────────────────────────────────────────────────────────────────────────

\subsection{Motivación para EDA}

La arquitectura orientada a eventos fue seleccionada por varias razones fundamentales:

\begin{infobox}[¿Por qué EDA?]
En un sistema donde un solo resultado de partido impacta en: puntajes de \textbf{N} usuarios,
generación de \textbf{N} notificaciones, y recalculo del ranking global — procesar todo de forma
síncrona generaría una latencia de respuesta proporcional a N. Con EDA, el administrador recibe
confirmación inmediata y el procesamiento ocurre en paralelo en background.
\end{infobox}

\subsection{Exchange y Routing Keys}

Se utiliza un \textbf{exchange de tipo topic} en RabbitMQ, que permite routing flexible basado
en patrones:

\begin{center}
\begin{tabular}{@{}lll@{}}
  \toprule
  \textbf{Evento} & \textbf{Routing Key} & \textbf{Productores} \\
  \midrule
  Resultado cargado    & \code{match.result.loaded} & API (Admin endpoint) \\
  Scores actualizados  & \code{scores.updated}      & Scoring Worker (encadenamiento) \\
  Jornada cerrada      & \code{round.closed}        & API (Admin endpoint) \\
  \bottomrule
\end{tabular}
\end{center}

\subsection{Flujo Completo: Carga de Resultado}

\begin{enumerate}
  \item El administrador POST a \code{/api/admin/matches/\{id\}/result}.
  \item El backend persiste el resultado en PostgreSQL y publica:
    \begin{lstlisting}[language=Python, caption=Publicación del evento]
await publisher.publish(
    routing_key="match.result.loaded",
    body={
        "match_id": match.id,
        "winner_player_id": winner,
        "tournament_id": match.tournament_id,
    }
)
    \end{lstlisting}
  \item RabbitMQ enruta el mensaje a \textbf{tres colas en paralelo} (fan-out mediante bindings):
    \begin{itemize}
      \item \textbf{Scoring Worker}: calcula puntos para cada predicción del partido.
        Publica \event{scores.updated} al finalizar.
      \item \textbf{Notification Worker}: genera notificación para cada jugador que predijo.
      \item \textbf{(inicio del chain)} El Scoring Worker encadena hacia Ranking.
    \end{itemize}
  \item El \textbf{Ranking Worker} consume \event{scores.updated} y recalcula posiciones globales.
\end{enumerate}

\subsection{Encadenamiento de Eventos (Event Chaining)}

Un patrón clave en el diseño es el encadenamiento: el Scoring Worker no solo actualiza puntajes,
sino que al completar su trabajo publica un nuevo evento \event{scores.updated}. Esto permite
que el Ranking Worker sea un consumidor desacoplado que no necesita saber nada de partidos o
predicciones — solo reacciona cuando los scores cambian.

Este patrón es análogo al patrón \textbf{Saga} en microservicios.

% ──────────────────────────────────────────────────────────────────────────────
\section{Stack Tecnológico}
% ──────────────────────────────────────────────────────────────────────────────

\subsection{Backend — FastAPI}

\begin{tabularx}{\linewidth}{@{}lX@{}}
  \toprule
  \textbf{Tecnología} & \textbf{Justificación} \\
  \midrule
  FastAPI 0.111   & Framework async nativo en Python, tipado via Pydantic, documentación OpenAPI automática. \\
  SQLAlchemy 2.0  & ORM moderno con soporte async, migrations, y type-safe queries. \\
  Pydantic 2.7    & Validación de datos en entrada/salida, serialización automática. \\
  Uvicorn         & Servidor ASGI de alto rendimiento basado en uvloop. \\
  aio-pika 9.4    & Cliente RabbitMQ async para Python, soporta reconexión automática. \\
  \bottomrule
\end{tabularx}

\vspace{0.5cm}

La API expone 6 routers RESTful:

\begin{itemize}
  \item \code{/api/auth} — Login y listado de usuarios
  \item \code{/api/rounds} — Jornadas y partidos
  \item \code{/api/predictions} — CRUD de pronósticos
  \item \code{/api/admin} — Carga de resultados, cierre de jornadas
  \item \code{/api/ranking} — Tabla de posiciones global
  \item \code{/api/notifications} — Alertas por usuario
\end{itemize}

\subsection{Frontend — React + TypeScript}

\begin{tabularx}{\linewidth}{@{}lX@{}}
  \toprule
  \textbf{Tecnología} & \textbf{Justificación} \\
  \midrule
  React 18.3       & Biblioteca UI con virtual DOM, hooks, componentes reutilizables. \\
  TypeScript 5.4   & Tipado estático end-to-end, mejor DX y refactoring seguro. \\
  Vite 5.2         & Build tool ultra-rápido con HMR, optimización de bundles. \\
  React Router 6   & Routing client-side declarativo. \\
  Tailwind CSS 3   & Utility-first CSS con design tokens personalizados. \\
  Material Symbols & Iconografía consistente con sistema de diseño Google Material 3. \\
  \bottomrule
\end{tabularx}

\subsubsection{Diseño Visual — Sistema de Tokens}

El frontend implementa un sistema de diseño oscuro inspirado en Material You (Material Design 3)
con los siguientes tokens de color primarios:

\begin{center}
\begin{tabular}{@{}lll@{}}
  \toprule
  \textbf{Token} & \textbf{Valor Hex} & \textbf{Uso} \\
  \midrule
  \code{surface}             & \#101415 & Fondo principal \\
  \code{surface-container}   & \#1d2022 & Cards y paneles \\
  \code{orange-accent}       & \#FF5C00 & CTAs, highlights, énfasis \\
  \code{secondary-container} & \#ff5e07 & Nav activo, botones primarios \\
  \code{on-surface}          & \#e0e3e5 & Texto principal \\
  \code{outline-variant}     & \#44474e & Bordes y separadores \\
  \bottomrule
\end{tabular}
\end{center}

\subsection{Broker de Mensajería — RabbitMQ}

RabbitMQ fue seleccionado sobre alternativas como Kafka o Redis Pub/Sub por:

\begin{itemize}
  \item \textbf{Persistencia de mensajes}: garantía de entrega incluso con workers caídos.
  \item \textbf{Exchange tipo topic}: permite routing flexible con routing keys jerárquicas.
  \item \textbf{Management UI}: interfaz visual para monitoreo de colas y mensajes.
  \item \textbf{Mature ecosystem}: cliente \code{aio-pika} con soporte async nativo.
  \item \textbf{Reconexión automática}: el patrón base en \code{workers/base.py} maneja
        desconexiones transparentemente.
\end{itemize}

\subsection{Base de Datos — PostgreSQL}

\begin{tabularx}{\linewidth}{@{}lX@{}}
  \toprule
  \textbf{Tabla} & \textbf{Descripción} \\
  \midrule
  \code{users}         & Jugadores y administradores \\
  \code{tournaments}   & Torneos ATP (ej: US Open, Wimbledon) \\
  \code{rounds}        & Jornadas de predicción por torneo \\
  \code{matches}       & Partidos individuales, con flag \code{is\_final} \\
  \code{predictions}   & Pronósticos usuario-partido con puntos calculados \\
  \code{rankings}      & Posición y puntos acumulados por usuario \\
  \code{notifications} & Alertas por usuario con estado de lectura \\
  \bottomrule
\end{tabularx}

\subsection{Infraestructura — Docker Compose}

El sistema completo se orquesta con Docker Compose en 5 servicios:

\begin{lstlisting}[language=bash, caption=Servicios Docker Compose]
services:
  frontend    # React/Nginx  - puerto 80
  backend     # FastAPI      - puerto 8000
  workers     # Python workers (scoring + notif + ranking + scheduler)
  postgres    # PostgreSQL   - puerto 5432
  rabbitmq    # RabbitMQ     - AMQP 5672, Management UI 15672
\end{lstlisting}

% ──────────────────────────────────────────────────────────────────────────────
\section{Modelo de Datos}
% ──────────────────────────────────────────────────────────────────────────────

\subsection{Entidades Principales}

\subsubsection{Predicciones y Scoring}

El modelo de puntuación es simple y auditable:

\begin{infobox}[Regla de Puntuación]
\begin{itemize}
  \item \textbf{Pronóstico correcto}: +10 puntos base.
  \item \textbf{Partido marcado como FINAL} (\code{is\_final = true}): +2 puntos adicionales.
  \item \textbf{Pronóstico incorrecto}: 0 puntos (sin penalización).
  \item El campo \code{predictions.points} es \code{NULL} hasta que el partido se resuelve.
\end{itemize}
\end{infobox}

\subsubsection{Estado de Partidos y Jornadas}

\begin{center}
\begin{tabular}{@{}ll@{}}
  \toprule
  \textbf{Entidad} & \textbf{Estados posibles} \\
  \midrule
  \code{Round.status}  & \code{open} $\to$ \code{closed} \\
  \code{Match.status}  & \code{pending} $\to$ \code{finished} \\
  \bottomrule
\end{tabular}
\end{center}

Los usuarios solo pueden cargar o modificar pronósticos cuando la jornada está \code{open}
y el partido está \code{pending}. Este invariante se valida en el backend.

\subsection{Flujo de Estados}

\begin{center}
\small
\textbf{Jornada:} \texttt{open} $\xrightarrow{\text{Admin cierra}}$ \texttt{closed}

\vspace{0.3cm}

\textbf{Partido:} \texttt{pending} $\xrightarrow{\text{Admin carga resultado}}$ \texttt{finished}

\vspace{0.3cm}

\textbf{Predicción:} \texttt{points = NULL} $\xrightarrow{\text{Scoring worker}}$ \texttt{points = N}
\end{center}

% ──────────────────────────────────────────────────────────────────────────────
\section{Workers Asíncronos}
% ──────────────────────────────────────────────────────────────────────────────

\subsection{Arquitectura de Workers}

Todos los workers heredan de una clase base \code{BaseWorker} definida en
\code{workers/base.py}, que encapsula:
\begin{itemize}
  \item Conexión y reconexión automática a RabbitMQ via \code{aio-pika}.
  \item Configuración del exchange y bindings.
  \item Manejo de errores y logging.
\end{itemize}

\subsection{Scoring Worker}

\textbf{Consume:} \event{match.result.loaded}

\textbf{Responsabilidades:}
\begin{enumerate}
  \item Obtiene todas las predicciones del partido resuelto.
  \item Para cada predicción, calcula los puntos según las reglas de negocio.
  \item Actualiza \code{predictions.points} en PostgreSQL.
  \item Publica \event{scores.updated} para triggear el recálculo de ranking.
\end{enumerate}

\subsection{Notification Worker}

\textbf{Consume:} \event{match.result.loaded}, \event{round.closed}

\textbf{Responsabilidades:}
\begin{enumerate}
  \item Por \event{match.result.loaded}: genera una notificación por cada usuario que predijo
        el partido, indicando si acertó o no.
  \item Por \event{round.closed}: genera una notificación global a todos los jugadores
        activos indicando el cierre de la jornada.
\end{enumerate}

\subsection{Ranking Worker}

\textbf{Consume:} \event{scores.updated}

\textbf{Responsabilidades:}
\begin{enumerate}
  \item Agrega puntos totales por usuario desde la tabla \code{predictions}.
  \item Calcula conteo de aciertos y total de predicciones evaluadas.
  \item Actualiza (UPSERT) la tabla \code{rankings} con posiciones recalculadas.
  \item Las posiciones se asignan por \code{RANK() OVER (ORDER BY total\_points DESC)}.
\end{enumerate}

\subsection{Scheduler}

El sistema incluye un \textbf{scheduler} (\code{workers/scheduler.py}) que monitorea
automáticamente el estado de las jornadas y puede cerrarlas cuando vence su deadline,
publicando el evento \event{round.closed} sin intervención manual del administrador.

% ──────────────────────────────────────────────────────────────────────────────
\section{Interfaz de Usuario}
% ──────────────────────────────────────────────────────────────────────────────

\subsection{Diseño Responsive}

El frontend implementa un layout adaptativo con dos modos:

\begin{description}
  \item[Mobile (\textless 1024px)] Top bar de 56px con logo + usuario, contenido en full-width,
        bottom nav redondeada de 70px con 4 tabs (Dashboard, Rankings, Alertas, Stats).
  \item[Desktop ($\geq$ 1024px)] Sidebar fija de 256px a la izquierda con navegación vertical,
        contenido principal con padding lateral de 32px.
\end{description}

\subsection{Vistas del Sistema}

\subsubsection{Login}
Selección de cuenta simplificada — múltiples usuarios pre-seeded para demostración.
Identifica automáticamente cuentas de administrador.

\subsubsection{Dashboard de Pronósticos}
Vista principal del jugador. Muestra:
\begin{itemize}
  \item Barra de progreso de pronósticos completados en la jornada activa.
  \item Cards de partidos con selector radio para elegir ganador.
  \item Indicadores visuales de partidos bloqueados, finalizados y puntos obtenidos.
\end{itemize}

\subsubsection{Rankings}
\begin{itemize}
  \item Podio visual para top 3 (oro, plata, bronce) con gradientes de color.
  \item Lista completa con posición, iniciales, aciertos y puntos.
  \item Auto-actualización cada 5 segundos (polling).
  \item Indicador de live con timestamp de última actualización.
\end{itemize}

\subsubsection{Estadísticas Personales}
\begin{itemize}
  \item Ring chart SVG de precisión general.
  \item Grid de 4 métricas: puntos totales, posición global, aciertos, evaluados.
  \item Historial completo de pronósticos con estado (pendiente/correcto/incorrecto).
\end{itemize}

\subsubsection{Notificaciones}
\begin{itemize}
  \item Lista cronológica de alertas por tipo (resultado, cierre de jornada).
  \item Iconografía codificada por color: verde (acierto), rojo (cierre), naranja (info).
  \item Marcado de leído al hacer clic. Badge de no leídos en el nav.
\end{itemize}

\subsubsection{Panel de Administración}
\begin{itemize}
  \item Métricas: jornadas abiertas y total de partidos pendientes.
  \item Tarjeta explicativa de la arquitectura de eventos (transparencia técnica).
  \item Por cada jornada: botón de cierre y cards por partido con botones de resultado.
  \item Al cargar resultado: dispara el pipeline de eventos completo.
\end{itemize}

% ──────────────────────────────────────────────────────────────────────────────
\section{Patrones y Decisiones de Diseño}
% ──────────────────────────────────────────────────────────────────────────────

\subsection{Patrón Publisher/Subscriber}

El sistema implementa Pub/Sub clásico donde:
\begin{itemize}
  \item El \textbf{publisher} (API) no sabe quiénes son los subscribers.
  \item Los \textbf{subscribers} (workers) no saben quién publicó el evento.
  \item El \textbf{broker} (RabbitMQ) actúa como intermediario desacoplado.
\end{itemize}

\subsection{Patrón CQRS Ligero}

Si bien no es CQRS estricto, el sistema separa semánticamente:
\begin{itemize}
  \item \textbf{Commands}: acciones que modifican estado (POST/PUT endpoints, events).
  \item \textbf{Queries}: lecturas de solo lectura (\code{GET /ranking}, \code{GET /predictions}).
\end{itemize}

\subsection{Idempotencia en Workers}

Los workers están diseñados para ser idempotentes: procesar el mismo evento dos veces
produce el mismo resultado que procesarlo una vez. Esto es crítico para garantizar
consistencia ante fallos y reenvíos del broker.

\subsection{Polling vs WebSocket}

La actualización del ranking utiliza polling cada 5 segundos desde el frontend.
Esta decisión fue deliberada:
\begin{itemize}
  \item \textbf{Simplicidad}: no requiere infraestructura de WebSocket.
  \item \textbf{Suficiencia}: el ranking no requiere actualización sub-segundo.
  \item \textbf{Correctness}: cada fetch garantiza datos frescos de PostgreSQL.
\end{itemize}

% ──────────────────────────────────────────────────────────────────────────────
\section{Deployment e Infraestructura}
% ──────────────────────────────────────────────────────────────────────────────

\subsection{Contenedores Docker}

Cada servicio tiene su propio \code{Dockerfile} optimizado:

\begin{itemize}
  \item \textbf{Frontend}: Multi-stage build — Node para compilar, Nginx Alpine para servir.
        El bundle resultante es de $\sim$196KB (60KB gzip).
  \item \textbf{Backend}: Python 3.x con dependencias en \code{requirements.txt}.
  \item \textbf{Workers}: Imagen compartida base con el backend, punto de entrada diferente.
  \item \textbf{PostgreSQL}: Imagen oficial con volumen persistente.
  \item \textbf{RabbitMQ}: Imagen con management plugin habilitado.
\end{itemize}

\subsection{Configuración Nginx}

El frontend incluye configuración Nginx (\code{nginx.conf}) con:
\begin{itemize}
  \item Serving de archivos estáticos con cache headers.
  \item Fallback a \code{index.html} para client-side routing (SPA mode).
  \item Gzip compression habilitada.
\end{itemize}

\subsection{CI/CD}

El repositorio incluye GitHub Actions configurados en \code{.github/} para
automatizar el pipeline de integración y deployment continuo.

% ──────────────────────────────────────────────────────────────────────────────
\section{Análisis de Calidad}
% ──────────────────────────────────────────────────────────────────────────────

\subsection{Puntos Fuertes}

\begin{enumerate}
  \item \textbf{Escalabilidad horizontal}: los workers son stateless y se pueden escalar
        simplemente aumentando instancias — RabbitMQ distribuye la carga automáticamente.
  \item \textbf{Resiliencia ante fallos}: si un worker cae, los mensajes permanecen en la
        cola y se procesan al reiniciar. No hay pérdida de datos.
  \item \textbf{Observabilidad}: la Management UI de RabbitMQ permite ver en tiempo real
        el estado de colas, mensajes pendientes y throughput.
  \item \textbf{Type safety end-to-end}: TypeScript en frontend + Pydantic en backend
        minimizan errores en tiempo de ejecución.
  \item \textbf{Developer Experience}: Vite HMR para desarrollo frontend, FastAPI auto-docs
        en \code{/docs} para testing de API.
\end{enumerate}

\subsection{Limitaciones y Trabajo Futuro}

\begin{tabularx}{\linewidth}{@{}lX@{}}
  \toprule
  \textbf{Limitación} & \textbf{Mejora Propuesta} \\
  \midrule
  Polling de ranking (5s)  & Migrar a WebSocket o SSE para updates push. \\
  Auth por selección de user & Implementar JWT con contraseñas hasheadas (bcrypt). \\
  Sin tests automatizados  & Agregar pytest para backend, Vitest para frontend. \\
  Workers en monolito      & Separar en microservicios independientes con su propio repositorio. \\
  Sin rate limiting        & Implementar límites por IP/usuario en la API. \\
  \bottomrule
\end{tabularx}

% ──────────────────────────────────────────────────────────────────────────────
\section{Conclusiones}
% ──────────────────────────────────────────────────────────────────────────────

Pick \& Serve demuestra en forma concreta cómo una \textbf{arquitectura orientada a eventos}
resuelve el problema del acoplamiento temporal en sistemas con múltiples efectos secundarios.

La elección de RabbitMQ como broker permitió implementar un fan-out genuino — un único
evento de resultado de partido dispara concurrentemente scoring, notificaciones y (vía
encadenamiento) actualización de ranking — sin que ninguno de estos procesos conozca
la existencia de los demás.

El stack tecnológico (FastAPI + React + PostgreSQL) representa una elección pragmática
que prioriza productividad sin sacrificar performance, y es representativo de lo que se
usa en la industria para sistemas SaaS de escala mediana.

Como prueba de concepto, el sistema cumple su objetivo pedagógico: evidenciar en código
ejecutable los beneficios del desacoplamiento, la asincronía y la comunicación via
mensajes en sistemas distribuidos modernos.

\vspace{1cm}
\begin{center}
{\color{accent}\rule{0.6\linewidth}{1.5pt}}\\[0.3cm]
{\small\color{codegray} Pick \& Serve · Ingeniería de Software II · 2026}
\end{center}

\end{document}
